\subsection{极大序域的特征刻画。欧拉-拉格朗日定理}

VI 第 25 页命题 6 所表达的性质刻画了极大序域。更精确地说：

定理 3（欧拉-拉格朗日）。--- 设 K 为一个序域。则以下三个性质等价：

\begin{enumerate}
\def\labelenumi{\alph{enumi})}
\item
  域 \(\mathbf{K}(i)\) 是代数闭的（其中 \(i\) 表示 \(-1\) 的一个平方根）。
\item
  序域 \(\mathbf{K}\) 是极大的。
\item
  \(\mathbf{K}\) 的每个正元素都是平方，且 \(\mathbf{K}[\mathbf{X}]\) 的每个奇数次多项式在 \(\mathbf{K}\) 中有一个根。
\end{enumerate}

显然 \(a)\) 蕴含 \(b)\) ：事实上，\(K\) 在同构意义下只有两个代数扩张，即域 \(K\) 其自身及 \(K(i)\) ，后者无法定序，因为 \(-1\) 是一个平方。

\(b)\) 蕴含 \(c)\) 这一事实无非就是 VI 第 25 页的命题 6。

我们尚需证明 \(c)\) 蕴含 \(a)\) 。这将由接下来的两个命题得出。

命题 7. --- 设 K 是一个序域，其中每个正元素都是平方。则 K(i) 的每个元素都是平方，且 K(i) 上每个次数为 2 的多项式在 K(i) 中都有根。

让我们首先证明第二个断言可以化归为第一个。可以把二次多项式 \(aX^{2} + bX + c\) ( \(a \neq 0\) )写成如下形式，通常称为标准三项式形式：

\[
a \left(\left(X + (b / 2 a)\right) ^ {2} - \left(b ^ {2} - 4 a c\right) / 4 a ^ {2}\right).
\]

如果 \(d\) 是 \((b^2 - 4ac) / 4a^2\) 的一个平方根，则 \(d - (b / 2a)\) 是所考虑的二次多项式的根。

现在我们证明每个元素 \(a + bi\) ( \(a \in \mathbf{K}\) , \(b \in \mathbf{K}\) ) 都是平方；我们要寻找一个元素 \(x + yi\) 使得

\[
(x + y i) ^ {2} = a + b i;
\]

这转化为 \(x^{2} - y^{2} = a\) 与 \(2xy = b\) ，由此我们推断出

\[
(x ^ {2} + y ^ {2}) ^ {2} = a ^ {2} + b ^ {2}.
\]

设 \(c\) 表示 \(a^2 + b^2\) 的正平方根；则 \(c \geqslant |a|\) , \(c \geqslant |b|\) 且 \(x^2 + y^2 = c\) ，由此得出 \(x^2 = (c + a)/2\) 与 \(y^2 = (c - a)/2\) 。由于 \(c \geqslant |a|\) ，这些方程在 K 中可解，且若 \(x_0\) 与 \(y_0\) 是两个解，则 \(x_0^2 - y_0^2 = a\) 且 \(2x_0y_0 = \pm b\) ，我们通过取 \(x = x_0\) 与 \(y = b/2x_0\) 得到所需的平方根。

命题 8. --- 设 K 是一个交换域（特征任意），并设 \(\mathbf{K}'\) 是多项式 \(\mathbf{X}^2 + 1 \in \mathbf{K}[\mathbf{X}]\) 的一个分裂域（V，第 21 页）。假设：

\begin{enumerate}
\def\labelenumi{\alph{enumi})}
\item
  K{[}X{]} 中每个奇数次的多项式在 K' 中有根；
\item
  \(\mathbf{K}'[\mathbf{X}]\) 中每个次数为 2 的多项式在 \(\mathbf{K}'\) 中有根。则 \(\mathbf{K}'\) 是代数闭的。
\end{enumerate}

首先注意，只需证明 K{[}X{]} 中每个非常数多项式在 K' 中有根：若 \(K' = K\) 这是清楚的；若 \(K' \neq K\) ，则 \([K': K] = 2\) ；设 \(a \mapsto \bar{a}\) 表示 \(K'\) 的唯一异于恒等映射的 K-自同构；若 \(f \in K'[X]\) 且 \(\bar{f}\) 表示对 f 的系数应用 \(a \mapsto \bar{a}\) 所得到的多项式，则 \(f\bar{f} \in K[X]\) ；若 \(a \in K'\) 是 \(f\bar{f}\) 的根，则 a 或是 f 的根，或是 \(\bar{f}\) 的根；因此 a 或 \(\bar{a}\) 是 f 的根。

因此设 f 是 K 上的一个多项式，次数为 \(2^{n}p\) ，p 为奇数。我们将对 n 进行归纳，该性质在 n = 0 时由假设 a) 成立。设 E 是 K 的一个扩域，在其中 f 分裂为线性因子：

\[
f (\mathbf {X}) = \prod_ {i} \left(\mathbf {X} - a _ {i}\right).
\]

设 \(b \in K\) ；令 \(y_{ij} = a_i + a_j + ba_i a_j \in E\) 且

\[
h (\mathbf {X}) = \prod_ {i <   j} \left(\mathbf {X} - y _ {i j}\right) \in \mathrm{E} [ \mathbf {X} ].
\]

该多项式的系数是这些 \(a_i\) 的、系数在 \(\mathbf{K}\) 中的对称函数；因此它属于 \(\mathbf{K}[\mathbf{X}]\) （IV，第 62 页，定理 1）；由于它的次数为 \(2^n p(2^n p - 1)/2 = 2^{n-1}p'\) ( \(p'\) 奇数)，根据归纳假设它有一个根 \(y_{ij}\) 属于 \(\mathbf{K}'\) 。如果我们注意到这对所有 \(b \in \mathbf{K}\) 成立，并且 \(\mathbf{K}\) 是无限域（事实上，一个有限域，由于它具有任意大奇数次数的单生成扩张（V，第 94 页，命题 3），不能满足 \(a\) )），那么我们可以推断出至少存在一对 \((i,j)\) 使得

\[
a _ {i} + a _ {j} + b a _ {i} a _ {j} \in \mathbf {K} ^ {\prime} \quad \text { and } \quad a _ {i} + a _ {j} + b ^ {\prime} a _ {i} a _ {j} \in \mathbf {K} ^ {\prime},
\]

，其中 \(b \neq b'\) 。则 \(a_i + a_j\) 与 \(a_i a_j\) 是 \(K'\) 的元素，因此 \(a_i\) 与 \(a_j\) 也是，因为它们是二次方程

\[
x ^ {2} - \left(a _ {i} + a _ {j}\right) x + a _ {i} a _ {j} = 0.\tag{Q.E.D.}
\]

的根。关于命题 8 的推广以及基于伽罗瓦理论的另一种证明，参见 VI，第 46 页，习题 33。

设 K 是一个序域，并设 \(\mathbf{K}^{\prime} = \mathbf{K}(i)\) ；对于每个元素 \(z = a + bi\) 属于 \(K^{\prime}\) ，范数 \(z\overline{z} = a^{2} + b^{2}\) （相对于 K，III，第 544 页，例 1）是 K 的一个正元素，且仅当 z = 0 时为零。如果 K 中的每个正元素都是平方（特别地，若 K 是极大序域），则范数 \(z\overline{z}\) 的正平方根称为 z 的绝对值，记作 \(|z|\) 。由于 \(|zz^{\prime}|^{2} = |z|^{2} |z^{\prime}|^{2}\) ，我们有 \(|zz^{\prime}| = |z| \cdot |z^{\prime}|\) 。

此外，三角不等式

\[
\left| z + z ^ {\prime} \right| \leqslant \left| z \right| + \left| z ^ {\prime} \right|
\]

对每一对元素 \(z, z'\) 属于 \(\mathbf{K}'\) 都成立。事实上，如果 \(z = a + bi\) 且 \(z' = a' + b'i\) ，那么这个不等式等价于

\[
(a + a ^ {\prime}) ^ {2} + (b + b ^ {\prime}) ^ {2} \leqslant a ^ {2} + b ^ {2} + a ^ {\prime 2} + b ^ {\prime 2} + 2 \sqrt {(a ^ {2} + b ^ {2}) (a ^ {\prime 2} + b ^ {\prime 2})}
\]

因而也等价于

\[
(a a ^ {\prime} + b b ^ {\prime}) ^ {2} \leqslant (a ^ {2} + b ^ {2}) (a ^ {\prime 2} + b ^ {\prime 2})
\]

它可以写成 \((ab' - ba')^2 \geqslant 0\) 。

定理 3 使我们能够确定一个极大序域上的所有不可约多项式：

命题 9. —— 若 K 是一个极大序域，则 K{[}X{]} 中仅有的不可约多项式是一次多项式，以及二次多项式 \(aX^{2} + bX + c\) ，其中 \(b^{2} - 4ac < 0\) 。

由于 \(\mathbf{K}(i)\) 是代数闭的，K 的每个代数扩张、因而 K 上的每个不可约多项式，其次数为 1 或 2。为了看出哪些二次多项式是不可约的，只需考虑标准形式 \(a((\mathbf{X} + (b/2a))^2 - (b^2 - 4ac)/4a^2)\) （参见 VI，第 26 页，命题 7）。

注：--- 转化为标准三项式形式可得到这一更强的结果：多项式 \(aX^{2} + bX + c\) 在给定的序域 K 上具有常号，当且仅当 \(b^{2} - 4ac < 0\) ，且此时多项式的符号与 a 的符号相同。

